\section{Review Protocol and Evidence Base}
\label{sec:methodology}

We use an adapted PRISMA-style reporting workflow; the protocol was not prospectively registered. The goal is to make the survey's evidence base transparent without letting the protocol dominate the story. We distinguish a \emph{Core Evidence Corpus}, containing attack studies with an explicit attacker, target, method, and outcome, from a \emph{Contextual Reference Corpus}, containing model substrates, benchmarks, adversarial-ML methods, robot-security work, and action-representation foundations. The frozen supplementary files provide the full record list, field dictionary, audit trail, and commands for regenerating the counts used in the main text.

\subsection{Databases, Dates, and Search Strings}

The last search for this revision was conducted on 16 June 2026. The search time range was all records indexed or publicly available by that date; no lower date bound was imposed, but direct VLA attack screening emphasized 2023--2026. We searched Google Scholar, Semantic Scholar, arXiv, IEEE Xplore, ACM Digital Library, USENIX proceedings, OpenReview, CVF Open Access, ACL Anthology, PMLR, RSS/CoRL/ICRA/IROS proceedings pages, and project/repository pages linked from papers. The supplementary README lists source-specific query templates, discovery-source codes, and commands for regenerating the statistics.

The core search strings were instantiated with spelling variants and venue filters:
\begin{itemize}
    \item \texttt{("vision-language-action" OR VLA OR "robot foundation model") AND (attack OR adversarial OR jailbreak OR backdoor OR poisoning OR privacy OR "membership inference")}
    \item \texttt{("OpenVLA" OR "RT-2" OR "Octo" OR "pi0" OR "LIBERO") AND (patch OR prompt OR backdoor OR "action token" OR "action chunk" OR trajectory)}
    \item \texttt{("embodied AI" OR "LLM-controlled robot" OR "robotic agent") AND ("prompt injection" OR jailbreak OR red-team OR attack)}
    \item \texttt{("robot" OR ROS OR "cyber-physical") AND (security OR sensor attack OR middleware attack OR HRI attack)}
    \item \texttt{("adversarial patch" OR "physical adversarial" OR LiDAR OR RAG OR "tool agent") AND (robot OR embodied OR VLA)}
\end{itemize}

Targeted verification searches were then run for direct VLA attack names and classes mentioned in recent papers: DropVLA/TabVLA, VLA-Fool, SABER, Q-DIG, BadVLA, GoBA, AttackVLA/BackdoorVLA, FreezeVLA, EDPA, ADVLA, UPA-RFAS, VLA-Hijack, Tex3D, TRAP, ReasonBreak, SilentDrift, trajectory redirection, partially observable patch attacks, and membership inference/VLALeaks.

Backward snowballing followed references from seed VLA attack papers, VLA/RFM surveys, the closest VLA safety survey, and embodied-security surveys when a cited work either attacked a relevant system or defined an attacked substrate. Forward snowballing used citation search from seed direct VLA attack papers and was stopped when new results no longer met inclusion criteria or were duplicates. Targeted verification search was used only after a named attack or project was encountered in a paper, review, project page, or reviewer suggestion; this reduces name-missing errors but is treated as a possible bias in Section~\ref{sec:open-problems}.

\subsection{Inclusion and Exclusion Criteria}

We included a work in the \emph{Core Evidence Corpus} if it satisfied at least one of the first two criteria: (i) it directly attacks a VLA model, reasoning-enabled VLA, VLA-controlled robot, VLA training channel, or VLA privacy channel; or (ii) it attacks an embodied LLM/VLM planner, robot policy, HRI channel, physical sensing path, or embodied agent with an explicit attacker and action/robot consequence. We included a work in the \emph{Contextual Reference Corpus} if it provided attack-evaluation methodology, a foundation model, benchmark, action representation, robot middleware/security substrate, or adversarial-ML mechanism necessary to interpret the core evidence.

We excluded works that only report ordinary robustness or distribution shift without an adversarial objective; text-only jailbreaks with no embodied, tool, or action boundary; image-classification attacks with no grounding, planning, action, or physical consequence relevance; broad AI safety position papers without attack or evaluation mechanisms; and records whose bibliographic identity could not be verified from an arXiv, publisher, OpenReview, project, or repository page. Duplicate records were merged when they referred to the same paper across arXiv, project page, repository, OpenReview, or conference page. When both preprint and venue versions existed, the peer-reviewed venue version was used as the bibliographic record and the arXiv/project page was retained only for version or artifact notes.

\begin{table}[t]
\centering
\footnotesize
\caption{Transparent screening flow for frozen database \texttt{VLA-Attack-ER-v2026-06-16}. Counts describe this manuscript's adapted PRISMA-style reporting workflow; the protocol was not prospectively registered.}
\label{tab:review-flow}
\renewcommand{\arraystretch}{1.08}
\setlength{\tabcolsep}{3pt}
\begin{tabularx}{\columnwidth}{@{}p{0.48\columnwidth}rY@{}}
\toprule
\textbf{Screening stage} & \textbf{Count} & \textbf{Notes} \\
\midrule
Records captured before de-duplication & 276 & Database queries, arXiv/OpenReview/CVF pages, seed surveys, backward/forward snowballing. \\
Duplicate or superseded records removed & 51 & Same work across arXiv, project page, repository, conference version, or renamed preprint. \\
Title/abstract screened & 225 & Screened for VLA, embodied robotics, adversarial objective, privacy objective, or evaluation relevance. \\
Excluded at title/abstract & 113 & Ordinary robustness only, no embodied implication, text-only safety, or unrelated robotics. \\
Full texts assessed & 112 & Full paper, arXiv page, publisher page, OpenReview page, project page, or repository was checked. \\
Full texts excluded & 27 & No adversarial objective (8); no action/robot boundary (6); insufficient VLA relevance (7); unverifiable identity (4); short duplicate (2). \\
Core Evidence Corpus & 30 & Direct VLA attacks (23) plus embodied-adjacent attacks with explicit attacker, target, method, and outcome (7). \\
Contextual Reference Corpus & 55 & VLA/RFM substrates, benchmarks, adversarial-ML methods, RAG/security-agent context, ROS/cyber-physical security, and action-representation foundations. \\
Included records in frozen database & 85 & Core evidence (30) plus contextual references (55). \\
\bottomrule
\end{tabularx}
\end{table}


\subsection{Evidence Levels}

Every included record is labeled by corpus type and evidence stratum. The \emph{Core Evidence Corpus} contains direct VLA attacks and embodied-adjacent attacks with explicit attacker, target, method, and outcome. The \emph{Contextual Reference Corpus} contains analogical and methodological records used to define substrates, mechanisms, benchmarks, or evaluation concepts. \emph{Direct evidence} means the victim is a VLA model or a VLA-controlled robot system. \emph{Embodied-adjacent evidence} means the victim is an embodied LLM/VLM planner, a traditional robot policy, a physical sensing path, or an HRI system with physical or action consequences but not a VLA policy. \emph{Contextual/analogical evidence} provides mechanisms or evaluation lessons from classifiers, autonomous-driving perception, RAG, software agents, general RL, ROS, or cyber-physical systems. Cross-study findings about attacks use the core corpus as the denominator unless a direct-VLA subset is explicitly named.

We used the following decision rule for direct evidence. A record is coded as \emph{direct VLA} only if (1) the target is a named VLA model, a reasoning-enabled VLA, a VLA-controlled robot, or a VLA training/privacy channel; and (2) the measured outcome concerns VLA actions, trajectories, action representations, reasoning-to-action behavior, triggered VLA behavior, or VLA training-data membership. If the target is an embodied LLM/VLM planner, robot middleware, HRI channel, or sensor stack without a VLA policy measurement, the record is \emph{embodied-adjacent}. If the work provides an attack mechanism, benchmark substrate, or representation baseline without embodied VLA validation, it is \emph{analogical/methodological}. Ambiguous records are conservatively downgraded rather than upgraded.

\subsection{Coding Reliability and Audit Procedure}

The 23 direct VLA records were independently re-coded in a second audit pass for principal fields: evidence stratum, family, lifecycle, knowledge level, access channel, privilege level, highest validation tier, target representation, objective, action representation, real-robot actuation, transfer reporting, query-budget reporting, and recovery reporting. The second coder was given the frozen CSV and coding rules but not the manuscript's narrative claims. Initial agreement was 23/23 for evidence stratum, family, lifecycle, knowledge/access coding, highest validation tier, target representation, objective, and action representation. Initial agreement for real-robot actuation was 21/23; reconciliation changed RoboGCG and VLALeaks from ambiguous/reported actuation to no real-robot actuation, while FreezeVLA was kept conservative as NR. Reporting counts were regenerated from the reconciled CSV.

We also audited a stratified non-direct sample of 18 records: 10 embodied-adjacent and 8 contextual/analogical. Agreement was 17/18 for strict core/context placement and 18/18 for lifecycle, access channel, highest validation tier, target representation, and objective under metadata-only review. The one flagged case was software-agent evidence, which is retained as contextual unless it has robot or embodied action consequence. We report percent agreement rather than Cohen's $\kappa$ because the audited sample is small and several fields are multi-label with many NR/NA values. For multi-label fields, disagreements were resolved by conservative downgrading or NR. The direct-audit file and README are supplied with the database supplement.

\subsection{Supplementary Evidence Files}

The supplement contains five files: the per-paper evidence table, a corpus/source table, a direct-VLA audit table, a direct-VLA validation-flags table, and a README with field dictionary and commands for regenerating the main statistics. The most important rule is orthogonality: the same paper can have a vision entry surface, an action-token target representation, an integrity objective, black-box knowledge level, physical access channel, simulator validation, and task-failure consequence. Those are different axes, not competing categories.

The structured evidence base is used in three ways. First, it prevents literature inflation: direct VLA attacks, embodied-adjacent attacks, and contextual references are not treated as the same kind of evidence. Second, it prevents metric inflation: ASR values are not pooled when they refer to different objectives, denominators, validation settings, or robot consequences. Third, it makes gaps visible: a paper can be strong direct evidence for simulator action deviation and still weak evidence for physical transfer, recovery, or adaptive attacks against runtime monitors.
